\section{Agentes de IA no desenvolvimento}

A seção anterior mostrou que modernizar um sistema exige entender o que ele faz e manter suas regras. Um agente de IA para desenvolvimento, como definido na introdução, lê os arquivos do projeto, altera o código, executa comandos e roda testes. Esta seção mostra como isso funciona em uma modernização e como os estudos avaliaram esses agentes.

O estudo de \citeonline{alves2026legacy} mostra bem como o agente é preparado para a tarefa. Os autores usaram o Claude Code, versão 2.1.96, com o modelo Opus 4.6. Para cada funcionalidade, criaram uma pasta de trabalho separada, com três arquivos de instruções, que são textos onde se escrevem o objetivo, as regras e as convenções que o agente deve seguir. O código antigo ficava em uma pasta somente de leitura, e o projeto novo, em C\#, partia de um modelo pronto.

O pedido era sempre o mesmo, padronizado: migrar a funcionalidade da pasta, lendo as instruções primeiro e usando o código antigo só como referência. O agente trabalhou em uma única passada, e o primeiro autor supervisionou as sessões e verificou o resultado.

Outros estudos avaliam os agentes com \textit{benchmarks}, que são conjuntos de projetos e testes preparados para comparar agentes nas mesmas condições. \citeonline{may2025freshbrew} criaram o FreshBrew, com 228 projetos Java públicos que compilam com o Java 8, falham no Java 17 e têm alta cobertura de testes, isto é, uma parte grande do código é conferida por testes. \citeonline{liu2026migrationbench} criaram o MigrationBench, com 5.102 repositórios e um subconjunto de 300, também para a migração do Java 8 para o 17. \citeonline{li2026repomod} criaram o RepoMod-Bench, com 21 repositórios em 8 linguagens, cerca de 1,6 milhão de linhas de código e 11.616 testes, para conversões entre pares de linguagens, como de Python para Java.

Os \textit{benchmarks} permitem comparar agentes de forma justa, mas usam projetos públicos e preparados para teste. Isso mostra bem o desempenho do agente, mas não mostra como é usá-lo em um sistema real do dia a dia, ponto que volta na Seção~\ref{sec:desafios-literatura}. O processo da IBM, visto na seção anterior, completa o quadro: em vez de um agente que conduz tudo, a análise do código guia o modelo de linguagem. Nos dois casos, fica uma pergunta: quem decide se o código novo está certo? É o tema da próxima seção.
